\section{Introduction}\label{sec:intro}

Much of today's trading is done by programs that read recent prices, decide what to trade, and by trading move the prices they read. Trend-following funds, execution algorithms, market-making programs and, more recently, AI trading agents all work this way. When such a market misbehaves, supervisors ask who was trading and how fast those traders react. This paper asks a prior question: what does a market's price record reveal about its population of feedback traders, and for which decisions is that enough?

\paragraph{The answer in one sentence.} Prices reveal the market's \emph{feedback kernel}: the response, $m$ periods later, of total feedback flow to a unit move in the price signal. That kernel is all that prices identify, and it is all that anyone trading on public information needs. Group composition, meaning who the groups are and how strongly each trades, matters only for actions aimed at a particular group, and learning it from prices has a known and steep price.

\paragraph{The model.} The market has a price and $K$ groups of traders. Each group keeps an exponentially weighted average of past prices, with memory rate $r_i$, and trades in proportion to it with weight $w_i$. A desk that tracks a trend with a Kalman filter, or that learns by regularized online gradient descent, behaves exactly this way (Lemma~\ref{lem:micro}). Total feedback flow is then a convolution of past prices with the kernel
\[
 k_m=\sum_i w_i r_i^{\,m},\qquad m=0,1,2,\ldots,
\]
the moments of the \emph{merged measure} $M$, which puts weight $w_i$ at rate $r_i$. The paper rests on three results about this one object.

\begin{enumerate}[label=(\arabic*),leftmargin=*,itemsep=2pt]
\item \emph{The kernel is sufficient} (Theorem~\ref{thm:eqinv} and Proposition~\ref{prop:anylinear}). Add a value, noise traders, an insider and a dealer, each following an arbitrary causal rule, which covers Kyle, inventory-averse, learning and misspecified dealers. Two trader populations with the same kernel then produce the same price and order-flow paths, the same insider trades and profits and the same dealer inventories, realization by realization. The two games have the same strategy sets and payoffs, and hence the same equilibria. The proof needs only that aggregate feedback flow is a fixed linear function of past public signals, so the result holds for any linear trader population, not only exponential memories.
\item \emph{The kernel is all that prices identify} (Theorem~\ref{thm:merged}). Two populations give identical price responses to every trade sequence if and only if their merged measures coincide. Groups with the same memory cannot be separated however long the record. Groups with different memories can be separated in principle.
\item \emph{Learning more has a known price} (Theorems~\ref{thm:kcohort}, \ref{prop:esprit} and~\ref{prop:passive}). To separate $K$ groups whose memories lie within $\Delta\eta$ of each other takes of order $(\Delta\eta)^{-(4K-2)}$ experiments, and Prony's method attains this order. Waiting instead of trading does not help: passive observation needs at least the same order of observations.
\end{enumerate}

\paragraph{The results survive equilibrium.} In the base model the price impact and the persistence of mispricing are primitives. Section~\ref{sec:equilibrium} derives both from competitive dealers who dislike holding inventory. The equilibrium price is then the same closed loop, path by path, with effective impact and persistence and with the memory rates unchanged (Theorem~\ref{thm:eqreduced}). So the kernel remains the sufficient statistic, and the cost of going beyond it keeps its order. The supplement adds a strategic insider with short-lived information. It also shows that, with positive feedback, liquidity disappears beyond a threshold in dealer risk aversion that falls with the feedback intensity.

\paragraph{What the kernel misses: limits.} Real feedback traders have position limits, margin and risk budgets. Section~\ref{sec:saturating} gives a two-layer theory for traders whose positions saturate.
\begin{itemize}[itemsep=1pt,topsep=2pt]
\item The sufficient object becomes a capacity-weighted measure over memory and leverage, of which the kernel is one linear functional (Theorem~\ref{thm:satsuff}).
\item Calm-market prices reveal the kernel and nothing more. Under hard limits, two populations with the same kernel give identical prices until some trader reaches its limit (Theorem~\ref{thm:stress}). Calm data therefore carry no information about how close traders are to their limits, and Gaussian data carry exponentially little. Under smooth limits, the information grows like the sixth power of the probe (Proposition~\ref{prop:smooth}).
\item Stress reveals everything that matters. Responses to probes of all sizes identify the net capacity at every memory and leverage level, which is exactly the equivalence class (Theorem~\ref{thm:satident}).
\item The results extend to traders who deleverage or liquidate outside their bands (Proposition~\ref{prop:delever}). An informed dealer who prices next-period flow leaks leverage information into calm prices, but only at the scale of the probability of a breach (Proposition~\ref{prop:dealers}).
\end{itemize}

\paragraph{Supervision splits in two.} A market-wide action, such as a uniform cut in activity to meet a ceiling on the feedback budget, depends on the population only through the kernel. A targeted action, such as halting the largest group, does not. Two populations whose price responses differ only by order $(\Delta\eta)^3$ can need different halts (Proposition~\ref{prop:target}). For such decisions, flow reported separately by trader class turns each group into an ordinary regression, with no penalty as the groups' memories approach each other (Proposition~\ref{prop:tagged}). It cuts the number of experiments needed by a factor of order $(\Delta\eta)^{-(4K-2)}$ (Theorem~\ref{thm:discvalue}).

\subsection{What is known and what is new}\label{sec:lit}

The central idea of this paper, that an aggregate is determined by a distribution of individual coefficients and not by the individuals, is old. We state its precedents first, because they fix what we can claim.

\paragraph{Aggregation of dynamic models.} \citet{granger1980} showed that aggregating many first-order autoregressions with different coefficients produces a process whose dynamics are governed by the distribution of the coefficients; with enough mass near one the aggregate has long memory. \citet{zaffaroni2004} develops the theory for large economies. Recovering the coefficient distribution from aggregate data is known as the disaggregation problem \citep{celov2008}. Our merged measure $M$ plays the role of that coefficient distribution, and Theorem~\ref{thm:merged}(i) is the corresponding identification statement for a closed feedback loop. \citet{forni1997} develop the general theory of what aggregate dynamics reveal about heterogeneous individual dynamics, and \citet{pesaran2014} the large-panel version.

\paragraph{Sufficient statistics for aggregate dynamics.} In macroeconomics, \citet{auclert2024} show that the matrix of intertemporal marginal propensities to consume is a sufficient statistic for the output response to fiscal policy, whatever the household heterogeneity behind it. The kernel $k$ plays the same role for a market with feedback traders. It is the response of aggregate demand to past prices, and it summarizes every heterogeneity that matters for prices.

\paragraph{Observational equivalence.} \citet{sargent1976} showed that two macroeconomic theories with opposite policy implications generate the same reduced form, so that data from one policy regime cannot decide between them. Theorem~\ref{thm:eqinv} is a result of this kind for trader populations, and Section~\ref{sec:supervision} draws the policy consequence: which supervisory actions the equivalence class supports, and which need data from outside it.

\paragraph{Mixture estimation.} The order $(\Delta\eta)^{-(4K-2)}$ is the known minimax rate for finite mixtures with $K$ near-colliding components \citep{heinrich2018,wu2020}, and super-resolution theory gives the same conditioning for clustered exponentials \citep{moitra2015,batenkov2021,katz2023}. We do not claim the exponent.

\paragraph{Heterogeneous horizons under impact.} \citet{leclere2026} study a linear mean-field model in which agents with heterogeneous forecast horizons trade under market impact, and in which impact depends on aggregate positions rather than on trader identities. Their object is equilibrium and price regularity, not what an observer can identify or what supervision needs.

\paragraph{Feedback trading and insiders.} Positive-feedback traders and their effect on prices go back to \citet{delong1990}. \citet{barberis2015,barberis2018} model extrapolators with exponentially weighted memories, and \citet{brock1998} model heterogeneous rules. The insiders and dealers covered by Section~\ref{sec:suff} include those of \citet{kyle1985} and of inventory models in the tradition of \citet{hostoll1981} and \citet{grossman1988}.

\paragraph{What is new.} Five things, each of which uses the market setting.
\begin{enumerate}[label=(\alph*),leftmargin=*,itemsep=2pt]
\item \emph{Sufficiency inside a strategic market.} The aggregation literature concerns the statistical law of an aggregate series. Theorem~\ref{thm:eqinv} is a statement about games: it holds path by path for every insider strategy and every causal dealer rule, so it transfers to every equilibrium concept defined from payoffs and public histories. Part~(iv) shows where it stops, since group-level quantities are not invariant.
\item \emph{Memories are zeros of the closed loop.} In mixture and open-loop problems the unknown rates are locations or poles of the observed law. In a feedback market they are transmission zeros of the price response (Theorem~\ref{thm:minphase}). That decides which estimators are valid: fitting poles to prices finds the closed-loop eigenvalues and misses the memories. Because the closed loop is minimum phase, passive data face at least the same barrier as active probing.
\item \emph{The results hold in equilibrium.} With inventory-averse dealers who price the feedback traders' predictable flow, the equilibrium is still the closed loop, with unchanged memories (Theorem~\ref{thm:eqreduced}).
\item \emph{Limits are invisible until reached, and stress reveals them exactly.} With saturating traders, the sufficient object is a measure over memory and leverage. Calm-market prices reveal only its linear part, and stress identifies all of it (Section~\ref{sec:saturating}). The last step of the stress result rests on a known uniqueness theorem for networks with one hidden layer \citep{sussmann1992}; the reduction from prices to that setting is ours.
\item \emph{A supervisory split with a price.} The kernel suffices for market-wide actions and not for targeted ones, and class-tagged disclosure removes the cost of the targeted information (Propositions~\ref{prop:target} and~\ref{prop:tagged}, Theorem~\ref{thm:discvalue}). A group that can mimic another group's memory also caps what any price-based targeting can achieve, whatever the amount of data (Proposition~\ref{prop:hiding}).
\end{enumerate}
The induction behind Theorem~\ref{thm:eqinv}(iii) is short. Its content is not the induction but the object it isolates. The kernel is identified from prices, it is sufficient for every price-based outcome, and it is insufficient exactly for the targeted decisions supervisors care about.

\paragraph{Machine checking.} The core statements are checked in Lean~4 against Mathlib with standard axioms only. The checks cover the sufficiency theorem (pathwise, for arbitrary causal rules), the $K$-group lower and upper bounds, the minimum-phase identification, the finite-sample passive bound, the reduced form of the inventory equilibrium, and the reduced form and best response in the insider equilibrium. Each statement says what is checked and what is proved only on paper; Supplement~\ref{app:verification} lists the declarations.

\paragraph{Outline.} Section~\ref{sec:model} sets out the model. Section~\ref{sec:suff} proves that the kernel is sufficient and that it is what prices identify. Section~\ref{sec:equilibrium} shows that both hold in equilibrium. Section~\ref{sec:cost} prices the information beyond the kernel. Section~\ref{sec:supervision} shows which supervisory decisions need it and what disclosure buys. Section~\ref{sec:saturating} turns to saturating traders. Section~\ref{sec:scope} states the limits, and Section~\ref{sec:conclusion} concludes.
